\documentclass{article}        % 1. Document class (Preamble)
\usepackage{graphicx}          % 2. Packages (Preamble)
\usepackage{minted}

\title{Handwriting Model}
\author{By: Jia Jie}

\begin{document}               % <--- Body starts here
\maketitle

\tableofcontents
\newpage
\section {KNN Train Test Split}


\newpage
\section {Explaining Folds}
* Remember in 100\% of data, when do train split, it's usually 70\% 
data for training and 30\% testing.
\vspace{5px}

\begin{center} 
    In that 70\% of training data:
\end{center} 

\begin{center}
\begin{minted}{python}
    from sklearn.model_selection import KFold
    folds = KFold(n_splits = 5, shuffle = true, random_state = 42)
\end{minted}
\end{center}

How it works?
\begin{itemize}
    \item Let's say we have 700 pictures. We divide these into 140 pictures in each chunk.
    \item So we have 1,2,3,4,5 chunk.
    \item let k (not n splits) is 3, in this iteration, 1,2,3,4,5 is randomized
    \item For first round (1), 1,2,3,4 is Fold used for training, and 5 is used for validation
    \item Goes on until five rounds, validation data helps choose K
\end{itemize}
\begin{figure}[htbp]
    \centering % Centers the image on the page
    \includegraphics[scale=0.4]{KFold.png} % Insert file name without extension
    \label{fig:my_image} % Used to cross-reference this figure in text
\end{figure}

**Note: This is only setting the instructions of loading the data for one k-value!
\newpage

\section {Explaining GridSearch CV}
\begin{minted}{python}
from sklearn.model_selection import GridSearchCV
optimal_k = GridSearchCV(
    estimator=KNeighborsClassifier(),
    param_grid={"n_neighbors": range(1, 30)},
    cv=folds,
)
\end{minted}

\begin{itemize}
    \item The KFold is where each k is tested. 
    \item It runs each k, and each k runs KFold.
    \item Repeats until each iteration is done.
    \item Note that the code above hasn't load the data yet, it's just setting the instructions
\end{itemize}
\end{document}              
